% Version 3.6: result transport, canonical extraction and per-DATA acceptance.
\subsection{Published results and the global interpreter}
\label{sec:exports}
The Publish Layer carries complete regional results to Espresso.
The sender transports evidence; it does not choose the authoritative
global base or the grouping of atomic acceptance transactions.
An \tagg{EXPORT} retains its outer chain, version, expected-base and
interval fields for routing and diagnostics. For acceptance these are
\defterm{hints}. Signed result contents and authenticated source context
determine the chain, producing pair, ancestry and append permission.
A false hint neither authorizes invalid DATA nor vetoes a complete
valid result carried in the same record.

\paragraph{Supported result framing.}
The existing EXPORT body and dependency syntax has two finite sequences:
typed DATA objects and typed evidence objects. Each object has its own
decodable type and length boundary. Evidence includes signed proposals,
receipts, ancestry, incoming REGULAR or Force boundaries, historical
authorities and source proofs. A dependency can contain further framed
objects or name an exact public object as specified in
\cref{sec:record-completeness}. The decoder inventories all such objects,
including DATA and receipts supplied as ancestry; the sender's declared
output interval is not this inventory. A malformed object with a known
frame boundary is skipped without discarding its independently framed
neighbors. Object identities and dependency edges are recomputed from
the bytes, not trusted from a sender-supplied index.
Required native application context is part of the result's evidence.
Traverse each normalized object once; an unresolved cyclic dependency
does not form a complete proof. This traversal grants no validity to
the object merely because its framing was decoded.

Both the existing v3.5 EXPORT body/dependency form and this revision's
form use this interpretation. No new wrapper is required to use a
complete result already recorded with an incorrect base, outer version,
interval or ordering hint. The signed chain and protocol domain still
have to match this interpreter; merely changing an outer routing label
cannot change that domain. A record with no independently decodable
result supplies no candidate. This grammar does not purport to decode
an encrypted blob or infer the boundaries of arbitrary damaged bytes.

\paragraph{Result closures and DATA identity.}
A \defterm{result closure} consists of a direct FINAL receipt root and
the transitive objects needed to validate its DATA and ancestry.
Closures may share objects. A valid root also supplies covering evidence
for its same-session DATA ancestors under \cref{sec:coverage}; each
earlier session on the path still needs its own receipt. The normalized
identity of a DATA is
\[
 \operatorname{id}(B)=(\text{chain},\beta(B),r(B),\Hash(B)).
\]
Changing sender, receipt encoding or signer set does not change this
identity. A declared interval cannot hide a complete next DATA supplied
inside a closure. Conversely, an absent dependency does not become
present merely because a sender lists its hash.

The common completeness predicate is applied to the closure used for
the proposed DATA, not to unrelated batch items. Invalid extra DATA,
false hints and bogus dependencies outside that closure do not veto it.
Invalid material inside a required closure does invalidate that proof.
Try another complete proof of the same DATA when one proof fails; an
invalid first wrapper never creates a permanent duplicate marker.

\paragraph{Canonical extraction: $\tagg{DeriveNextDATA}$.}
Let $p=(q,n,j)$ identify an Espresso block, namespace transaction and
record. Process these records in that order after the block's prelude.
At the current state, inspect its normalized closures for DATA at
$h(G)+1$. Order candidate proof choices by the tuple
\[
 (\beta(B),r(B),\Hash(B),\Hash(D),\operatorname{pos}(\pi)),
\]
where $D$ is the direct-receipt root providing coverage and
$\operatorname{pos}(\pi)$ is the canonical frame path of the evidence.
Hashes use their fixed byte order. Equal normalized proof objects are
deduplicated. Ignore lower DATA heights except as evidence; never jump
to a higher height. Recompute the next-height candidates after a commit.
This is a finite deterministic scan under the record's input-size limit.

The incoming boundary fixes the only session and region permitted to
supply that next DATA. Within that session, regional safety makes two
valid committed next DATA compatible and hence identical. The proof
ordering chooses evidence, not a fork. Neither a later session's
receipt nor a lower hash can override the fixed incoming boundary.

\paragraph{Acceptance checks.}
For a candidate $B$, derive its actual base from $G$ and require:
\begin{enumerate}[label=(\roman*),itemsep=2pt]
\item its signed domain is this chain's and its producing branch
version is current and open;
\item its incoming boundary, historical selection and authority assign
it to the current export duty; its authenticated DATA ancestry extends
the exact current $G$, including hashes, heights and original provenance;
\item its height is $h(G)+1$, its direct or same-session covering
receipt is valid, and every required admission, parent, quota, reference
and native application-validity check succeeds;
\item all required Espresso selection references strictly precede
the containing block $J_q$, with their historical snapshots verified;
\item if Force created its version in $J_{q_f}$, then $q>q_f$.
\end{enumerate}
No check compares a sender's expected-base hint with $G$. Compatibility
is checked against the actual signed path instead. A record overlapping
$G$ may therefore supply its next compatible suffix. DRAIN can occur in
supporting ancestry, but is never appended as application DATA.

\paragraph{Commit one DATA: $\tagg{CommitNextDATA}$.}
On successful checks, atomically append $B$, increment its producing
session's accepted count and its region's accepted service count, and
remove that region from $\mathcal O$ if present. Count this DATA once.
A continuing duty keeps its window and receives the new window count;
an empty owed set is not refilled merely because another DATA arrived.
If $B=T_{\beta,r}$, cancel that duty's window and install the pending
REGULAR boundary described below in the same transaction. Durably expose
the updated $G$ and accounting before attempting a later DATA or
successor operation. The already saved pre-payload $\chi_E(q)$ and
$\mathcal O_E(q)$ are unchanged.

Acceptance is atomic per DATA and its associated accounting. There is
no all-or-nothing outer EXPORT transaction. An invalid later candidate
cannot roll back an accepted prefix. Batching is an implementation
optimization only if it preserves these committed prefixes and does
not postpone their exposure behind an unresolved successor operation.

\paragraph{Commit terminal DATA before resolving its successor.}
\label{sec:terminal-commit}
The terminal transaction stores the retained terminal anchor and its
commitment evidence, source context $(i_b,\beta,r)$, unchanged beta,
target slot $r+1$, signed selection reference and the regular selection
lower bound. It marks the boundary \defterm{pending resolution}; no
successor or DATA window is fabricated. Its owed ledger is the prior
ledger after the producing region's discharge, without a new round.
This transaction requires the completed session's incoming boundary
and authority, not $\tagg{AUTH}_{r+1,0}$.

Only the subsequent $\tagg{ResolveBoundary}$ operation obtains that
target authority, runs $\tagg{PrepareOwed}(r+1)$ and verifies the
prescribed historical REGULAR selection. A missing authority pauses
this operation after the terminal transaction has committed. A proven
absence of a candidate also leaves the boundary unassigned, but is
different from missing evidence. Neither case reverses accepted DATA,
restores the completed window, or opens a new window.
With evidence complete and no candidate yet, mark this resolution
attempt complete and continue source replay; a later prelude resumes
the prescribed historical search for the still-pending boundary.

Resolution resumes at its original canonical transition exactly once.
It preserves the historical selector and the public opening timestamp
in \cref{sec:ready}; local retrieval time never starts a new grace
period. An already fixed speculative handoff is checked against this
same history. Once resolution succeeds, the same record may supply a
successor DATA if its own references precede $J_q$. Thus normal
multi-session records remain possible. If the terminal publication
block is the first possible selection position, successor DATA still
cannot be accepted in that block.

Input authentication is incremental at this dependency boundary. It
authenticates the complete namespace and controls, then the authorities
needed for the prelude and current DATA transaction. It does not eagerly
demand every authority used by later operations in the record. In
particular, a preliminary whole-record check cannot wait for
$\tagg{AUTH}_{r+1,0}$ before committing an otherwise valid terminal DATA.

\paragraph{Durable subpositions and replay.}
\label{sec:subpositions}
The public cursor is
\[
 \zeta=(q,n,j,\nu,\mathit{stage}),\qquad
 \mathit{stage}\in\{\tagg{prelude},\tagg{scan},\tagg{resolve},\tagg{done}\}.
\]
Here $(q,n,j)$ is the block and record position, and $\nu$ counts DATA
transactions committed from that record. The prelude uses the sentinel
$n=j=\nu=0$. The stages mean prelude, next-DATA scan, pending boundary
resolution, and record complete, respectively. The boundary record
retains its original canonical resolution position across later scans.
At a DATA commit it also binds the normalized record commitment and
the resulting pending-boundary state. The next-height rule and this
cursor fix the continuation without dependence on private caches.
An implementation may retain a larger validation index as auxiliary
state. Duplicate source delivery resumes the cursor; it never reruns
a committed count or prelude.

An unresolved canonical boundary pauses subsequent replay at that
subposition. Earlier DATA commits remain exposed and final even though
the entire record or block is not yet complete. Once evidence is
available, resume the same transition with its original timestamp.
No out-of-order application of later controls is needed. Two observers
at different subpositions hold prefix-comparable DATA, just as they do
at different complete block heights.

If no candidate passes and no genuine canonical boundary wait remains,
finish the record and continue. An independently incomplete candidate
does not suspend other complete closures. A record skipped for a genuine
gap gains no future append permission: the interpreter never revisits
a completed past record to accept new DATA retroactively. Its objects
may remain historical evidence and may be referenced by a later record.
Closed-version DATA never becomes appendable through either path.
An incomplete namespace download, unlike an incomplete candidate,
prevents processing the block at all.

\paragraph{Scope of the transport interpretation.}
\label{sec:transport-scope}
The frozen delivery assumption supplies a record carrying the complete
regional result in the chain's retrievable namespace. Here ``carrying''
uses the supported decodable body/dependency syntax above, including
the same syntax in an earlier bad-hint EXPORT. It does not mean that
the result is accepted, and no exact-base submission or corrected-retry
deadline is added. A complete terminal result does not require its
successor authority. Missing payload or required ancestry was already
excluded from complete results in \cref{sec:delivery}.
The frozen service guarantee remains transport only: it does not promise
that a carried result passes this chain's acceptance rules. A stale
candidate with incompatible signed ancestry, a closed version or a missing
required predecessor still cannot append. This revision changes the
interpreter's treatment of a sender's base hint alone; it does not
change either delivery guarantee or infer acceptance from visibility.

If ``already recorded'' is instead interpreted to discharge delivery
for opaque bytes, an undiscoverable unrelated namespace, or DATA plus
a receipt with required ancestry omitted, the frozen model does not
give bounded access to a decodable complete result. The progress
argument below is not established under that broader interpretation;
an added access promise would be needed. This limitation cannot be
removed by demanding a new wrapper or declaring a wrong-base complete
result unrecorded. For the stated supported syntax, decoding and
per-DATA acceptance are protocol operations on the already promised
bytes, not a stronger publication service.
